% =====================================================================
\cvsection{Research \& Professional Experience}
\cventry{2024--Present}{Graduate Research Assistant}{Health Informatics \& Data Science, LUC}{Chicago, IL}{}{
\begin{itemize}
    \item \textbf{Multimodal Clinical Modeling:} Built an end-to-end multimodal fusion pipeline on MIMIC-IV (30-day post-discharge mortality), comparing early, intermediate, and late fusion of structured EHR data and clinical notes (TabPFN, XGBoost, BioClinical-ModernBERT; LoRA SFT of Llama-3; TIES / DARE-TIES adapter merging).
    \item \textbf{Clinical NLP / LLMs:} Developed a hybrid dense + BM25 retrieval system with an LLM-verification stage for schema-guided clinical information extraction of synthetic nursing dictations (MEDIQA-SYNUR 2026 shared task, LREC/ClinicalNLP).
    \item \textbf{Prehospital Stroke Triage:} Developed and validated ML models on EMS data to predict stroke and severe-stroke events; led the workflow / implementation study and clinician co-design.
    \item \textbf{Implementation Science:} Conducted qualitative analysis (CFIR) and customer discovery (I-Corps) to connect model development with real EMS workflows; collaborated with paramedics and clinicians on paramedic-centric interface design.
\end{itemize}}

\cventry{2023--2024}{Data Scientist}{Health Informatics \& Data Science, Parkinson School, LUC}{Chicago, IL}{}{
Aug 2023 -- Aug 2024. Preceded by a Bioinformatics Internship (May--Aug 2023; MS capstone).
\begin{itemize}
    \item Conducted a retrospective analysis of outpatient EHR data to evaluate hypertension management across a community health network.
    \item Investigated the association between Social Determinants of Health (Area Deprivation Index) and therapeutic inertia using statistical modeling.
    \item Disseminated findings via a moderated digital poster at AHA Scientific Sessions and a peer-reviewed publication in the \textit{Journal of Human Hypertension} (Nature Portfolio).
\end{itemize}}

